\section{Conclusion}

This paper presents an integrated AI-based framework for MSME viability assessment that advances the state of practice across three dimensions: multi-class prediction, transparent explainability, and prescriptive intervention guidance.

The framework introduces a five-class health taxonomy (Critical, At-Risk, Stable, Growing, Thriving) that transcends the binary classification paradigm dominant in existing MSME and credit risk assessment literature. Trained on 899,164 U.S. SBA loan records using XGBoost~\cite{10} and LightGBM~\cite{11} as primary prediction engines, the system achieves overall accuracies of 92.78\% and 92.98\% respectively, with macro AUC exceeding 0.986, substantially outperforming both naive baselines (53.31\%) and traditional Logistic Regression (77.67\%) on the five-class task.

The integration of SHAP-based explainability~\cite{17} provides per-prediction feature attribution that renders each viability assessment transparent and auditable. Global SHAP analysis identifies the gross approved amount and loan term as the dominant viability predictors, providing actionable insight for lending policy.

The counterfactual recommendation engine, inspired by DiCE~\cite{20}, extends the system from diagnostic prediction to prescriptive decision support. The engine achieves 100\% feasibility across 200 evaluated instances, resolving 82\% of Critical-class transitions through single-feature changes with a mean generation latency of 2.03 milliseconds --- demonstrating both the effectiveness and the real-time operational viability of the prescriptive approach. The convergence between SHAP-identified risk factors and counterfactual recommendations confirms internal system coherence.

The complete framework is deployed as an enterprise-grade decoupled architecture with persistent audit logging, addressing the production deployment challenges documented in the machine learning systems literature. We note that the health taxonomy is grounded in historical loan repayment outcomes; validation against independent enterprise survival data remains an important direction for future work. By combining graduated risk stratification, transparent explainability, and actionable prescriptions within a single operational system, the proposed framework offers a methodological contribution toward more accurate, transparent, and actionable MSME viability assessment --- a domain where the gap between available tools and practitioner needs remains substantial~\cite{6},~\cite{7}.
